\documentclass{standalone}
\usepackage{circuitikz}
\usetikzlibrary{calc}
\definecolor{circuitnotemark}{HTML}{FE00FE}
\begin{document}
\begin{circuitikz}[american, line width=0.8pt]
\ctikzset{bipoles/length=1.2cm}
\foreach \x in {0,2,4} {\foreach \y in {0,-2,-4} {\fill[gray, opacity=0.35] (\x,\y) circle (1.1pt);}}
\node[gray, font=\scriptsize] at (0,0.84) {1};
\node[gray, font=\scriptsize] at (2,0.84) {2};
\node[gray, font=\scriptsize] at (4,0.84) {3};
\node[gray, font=\scriptsize, anchor=east] at (-0.84,0) {a};
\node[gray, font=\scriptsize, anchor=east] at (-0.84,-2) {b};
\node[gray, font=\scriptsize, anchor=east] at (-0.84,-4) {c};
\coordinate (a1) at (0,0);
\coordinate (a3) at (4,0);
\coordinate (c3) at (4,-4);
\coordinate (c1) at (0,-4);
\draw (a1) to[R, l_=$R_{1}$, i>^=$i_{1}$] (a3); % line 3
\draw (c3) to[R, l_=$R_{2}$, i>^=$i_{2}$] (c1); % line 4
\path (8,-4.974) rectangle (16.403,0.593); % line 6
\node[anchor=west, inner sep=0, circuitnotemark, font=\large] at (8,0) {X}; % line 6
\node[anchor=west, inner sep=0, circuitnotemark, font=\large] at (8,-0.487) {X}; % line 6
\node[anchor=west, inner sep=0, circuitnotemark, font=\large] at (8,-0.974) {X}; % line 6
\node[anchor=west, inner sep=0, circuitnotemark, font=\large] at (8,-1.46) {X}; % line 6
\node[anchor=west, inner sep=0, circuitnotemark, font=\large] at (8,-1.947) {X}; % line 6
\node[anchor=west, inner sep=0, circuitnotemark, font=\large] at (8,-2.434) {X}; % line 6
\node[anchor=west, inner sep=0, circuitnotemark, font=\large] at (8,-2.921) {X}; % line 6
\node[anchor=west, inner sep=0, circuitnotemark, font=\large] at (8,-3.408) {X}; % line 6
\node[anchor=west, inner sep=0, circuitnotemark, font=\large] at (8,-3.895) {X}; % line 6
\node[anchor=west, inner sep=0, circuitnotemark, font=\large] at (8,-4.381) {X}; % line 6
\node[anchor=south west, inner sep=2pt, circuitnotemark, font=\LARGE\bfseries] (circuittitle) at (current bounding box.north west) {X};
\path (circuittitle.south west) rectangle ++(10.058,0.853);
\end{circuitikz}
\end{document}
